\subsection{معرفی و انگیزه}
این مکانیزم بر مبنای مبانی نظری پیشرفته در ادبیات توجه و فشرده‌سازی \cite{eckart1936approximation} توسعه یافته است.

یکی از مهم‌ترین یافته‌ها در زمینه فشرده‌سازی مدل‌های زبانی بزرگ، پدیده مقادیر پرت فعال‌سازی (\lr{Activation Outliers}) است. تحقیقات اخیر نشان داده‌اند که در مدل‌های ترانسفورمر، بخش بسیار کوچکی از ویژگی‌ها (کانال‌های خاص) دارای انرژی و بزرگی قابل توجه بالایی هستند؛ این کانال‌های پرت اگرچه از نظر تعداد اندک هستند، اما نقشی حیاتی در توانایی استدلال و توجه شبکه ایفا می‌کنند و هرگونه تقریب با رتبه پایین بر روی آن‌ها به فروپاشی عملکرد مدل می‌انجامد.

الگوریتم \lr{OATS (Outlier-Aware Training \& Saliency-Based Sparsification with Low-Rank Matrix Decomposition)} با هدف جداسازی و مدیریت اختصاصی این دو ماهیت ساختاری توسعه یافته است. در این شیوه، ماتریس وزن پس از پیش‌شرط‌سازی مبتنی بر گشتاور دوم فعال‌سازی‌ها، به دو مؤلفه متعامد تجزیه می‌شود: یک مؤلفه کم‌مرتبه متراکم (\lr{Low-Rank Component}) برای جذب جهت‌های پرت پرانرژی، و یک مؤلفه تنک ساختاریافته (\lr{Sparse Residual Component}) با الگوی $2:4$ برای حفظ جزئیات باقیمانده بدون تحمیل حافظه اضافه.
